\scanpage{41}
$S$ schéma régulier connexe de \unsure{t.f.|v.f.} sur $\unsure{\mathbb Q|Q}$

$k$ clôture alg. de $\mathbb Q$ dans $\Gamma(S,\unsure{\mathscr O|\mathcal O}_S)$, de sorte que $k/\mathbb Q$ fini, $S/k$ géom. \unsure{conn|conn.}

$k\hookrightarrow\mathbb C$ une immersion

$\mathcal M$ catégorie de motifs (``iso'', lisses $\frac12$ simples) sur $S$, $\mathcal M$ de type fini.

L'immersion $k\hookrightarrow\mathbb C$ définit une classe d'isomorphie de foncteurs fibres\inserted{ val\,/\,$\mathbb Q$} sur $\mathcal M$, savoir les $T_{\xi}\deleted{\unsure{T|P}}:M\longrightarrow T_B(M_{\xi}),\quad \xi\in S(\mathbb C)$\inserted{$\underset{\deleted{\operatorname{Hom}_k(\operatorname{Spec}\mathbb C,S)}}{\Vert}$}.

Soit \deleted{$G$ un groupe alg. /$\mathbb Q$ (réductif) isomorphe aux $\underline{\operatorname{Aut}}_{\otimes}(T_{\xi})$\unsure{en|.}.} \inserted{$T$ un foncteur fibre de cette classe, et \deleted{\unsure{soit}}}
\[
G=\underline{\operatorname{Aut}}_{\otimes}(T),\quad\text{donc }\mathcal M\approxeq\operatorname{Rep}(G)
\]

Alors
\[
P=\underline{\operatorname{Isom}}_{\otimes}(T_S,T_{\mathrm{DR}})
\]
est un torseur sous $G_S$. Il est muni d'une \uline{connexion} intégrable relativement à $k$ (ou $\mathbb Q$, c'est kif kif) compatible avec $G$,

De plus le foncteur $T_{\mathrm{DR}}$, isomorphe par construction de $P$ à
\[
M\longmapsto P\mathbin{\times}^{G_S}M_S\qquad\text{sur }\operatorname{Rep}(G)
\]
est muni d'une $\otimes$-\uline{filtration}

\deleted{Soit $S_{\mathbb C}$} Considérons $S(\mathbb C)$, et le revêtement $\deleted{\bar S}\unsure{\widetilde P|\bar P}_{\mathbb C}$ de $S(\mathbb C)$ dont la fibre en $\xi\in S(\mathbb C)$ est formée des isom $T\longrightarrow T_{\xi}$. Alors $\bar S_{\mathbb C}$ est un revêtement principal de $S(\mathbb C)$, de groupe $G(\mathbb Q)$, \deleted{qui n'est autre que le rev. \unsure{associé}, \unsure{fibré} \unsure{pr.} assoc. à $P_{\mathbb C}^{\mathrm{an}}$ sur $S_{\mathbb C}^{\mathrm{an}}\underset{S(\mathbb C)}{\Vert}$} et \unsure{un|on} \deleted{\unsure{une}} \deleted{\unsure{avec|de} $S_{\mathbb C}$} \deleted{\unsure{application}} \inserted{morphisme (hol)} évident
\[
p:\deleted{\unsure{\bar S|\bar P}}\widetilde P_{\mathbb C}
\longrightarrow P_{\mathbb C}^{\mathrm{an}}\deleted{\unsure{G|g}}
\]
satisfaisant
\[
\left[\begin{aligned}
&p(\unsure{\xi|g}\,\times\,g)=p(x)g\\
&p\ \text{compatible aux connexions \unsure{sur un|sur nos|sur}}\\
&\hspace{4em}S\text{-fibrés}\ldots
\end{aligned}\right.
\]

\[
\deleted{f:\bar S_{\mathbb C}\longrightarrow\unsure{\gamma|\mathscr Y}_{G}}
\]
\deleted{où $\gamma_G$ est l'espace des structures de Hodge sur $G_{\mathbb R}$ (compatibles, si on y tient, avec un $i:\mathbb G_m\relbar G$ fixé). Alors $f$ satisfait aux conditions}

\deleted{a) $f$ holomorphe}

\deleted{b) $f(xg)=g^{-1}f(x)$ pour $g\in G(\mathbb Q)$ [\uline{NB} $G(\mathbb R)$ donc $G(\mathbb Q)$ opère sur $\gamma_G$ par autom. \unsure{intérieurs|intégrieurs},]}

\deleted{c) .}
% archive sheet number: 40
